\documentclass{article}

\usepackage{microtype}
\usepackage{booktabs}
\usepackage{hyperref}
\makeatletter
\def\input@path{{../attachments/icml2026/}}
\makeatother
\usepackage{icml2026}

\hypersetup{
  pdftitle={Evidence-Bound Review Under a Fixed Deadline: A Reliability Contract for ICML-Style Review},
  pdfauthor={Anonymous},
  pdfsubject={Anonymous Track 2 technical report},
  pdfkeywords={scientific review, evidence trace, reliable orchestration}
}
\icmltitlerunning{Evidence-Bound Review Under a Fixed Deadline}

% The event requires a fully anonymous report. Preserve the official ICML
% preliminary notice while omitting the style's affiliation and correspondence
% wrapper, which otherwise emits demonstration identity fields in blind mode.
\makeatletter
\renewcommand{\printAffiliationsAndNotice}[1]{%
  \global\icml@noticeprintedtrue
  {\let\thefootnote\relax
    \footnotetext{\hspace*{-\footnotesep}\Notice@String}}%
}
\makeatother

\begin{document}

\twocolumn[
  \icmltitle{Evidence-Bound Review Under a Fixed Deadline: \\
    A Reliability Contract for ICML-Style Review}
  \vskip 0.3in
]

\printAffiliationsAndNotice{}

\begin{abstract}
Time-bounded scientific review must combine evidence tracing with reliable
batch execution. We implement an anonymous Track 2 review system that keeps
platform side effects in a single root coordinator, delegates paper-local
drafting to three bounded workers, and validates a canonical review schema
before any post attempt. Across three deterministic 10-paper synthetic-mock
runs, the runtime completed 30/30 assigned reviews with 100\% schema validity
and zero duplicate post attempts or idempotency keys; a separate fault run
completed 10/10 with schema validity 10/10, no duplicates, one recovery each
for malformed JSON output, worker timeout, claim timeout, and post timeout, and
one in-process ledger reopen. In an actual two-process dry-run, the first
process exited 75 after four mock-verified papers and a fresh process exited 0
at 10/10 with schema validity 10/10 and no duplicates while preserving the
four-paper ledger prefix and completed manifest and outbox bytes. On frozen
seeded synthetic issues, the deterministic reviewer scored TP 20, FP 0, FN 0,
and F1 1.0, compared with TP 10, FP 0, FN 10, and F1 0.6667 for the naive
baseline. These measurements are mock-only and do not establish live-platform
throughput or review quality on unseen papers.
\end{abstract}

\section{Problem and Contract}

Track 2 requires an ICML-style review agent and a structured review of a frozen
paper. The live operating window is 25 minutes, so review quality and batch
reliability are coupled: a well-reasoned draft is not useful if it is lost,
duplicated, or posted without confirmation. The system therefore treats
evidence identity, review validity, and side-effect state as one end-to-end
contract.

The official upstream skill is pinned at commit
\texttt{a9f4f258...153473f}. Its Track 2-only path
requires no new training, GPU, experiment service, or training metric. It
requires frozen paper and evidence hashes, an independently frozen agent
definition, a structured result, and explicit reporting when supplied evidence
is insufficient. The local wrapper preserves the upstream
\textit{Contribution} assessment while retaining the event fields
\textit{Significance}, \textit{Originality}, and \textit{Comment}; it does not
convert one field into another.

\section{System Design}

\subsection{Single Owner for Side Effects}

A root coordinator is the only component permitted to discover assignments,
claim papers, inspect remote status, post reviews, or update the shared ledger.
Three persistent review workers receive one frozen paper manifest at a time and
return only a structured \texttt{ReviewDraft}. This boundary prevents worker
retries from becoming duplicate platform actions.

\begin{figure}[t]
  \centering
  \fbox{\parbox{0.92\columnwidth}{
    \centering
    Serialized platform lane: discover, claim, reconcile, post, verify
    \par\smallskip
    $\downarrow$ frozen manifest and bounded lease
    \par\smallskip
    Three review workers: paper-local evidence analysis only
    \par\smallskip
    $\downarrow$ \texttt{ReviewDraft}
    \par\smallskip
    Deterministic schema validation and atomic ledger transition
  }}
  \caption{Authority is narrow by construction. Workers cannot produce remote
  side effects or mutate shared state.}
  \label{fig:architecture}
\end{figure}

\subsection{Frozen Inputs and Evidence Trace}

Before a worker starts, the coordinator records the paper identifier, paper
file identity, SHA-256 and size, evidence file identities, hashes and sizes,
aggregate input and evidence hashes, agent version, prompt hash, and schema hash
in a per-paper manifest. The frozen corpus supplies the allowed file paths
separately, and Root grants the active lease separately from the manifest.
Unsafe paper identifiers are rejected before any artifact or ledger write. A
worker must return the frozen identities and its Root-issued lease identity
with page, section, table, figure, or saved-result locations for central
strengths and weaknesses. An identity mismatch blocks execution and cannot
enter the targeted repair path.

\subsection{Canonical Review Shape}

The draft contains Summary, Strengths, Weaknesses, Questions, Ethics and
Limitations, Evidence Trace, and score rationales. It carries Soundness,
Presentation, Contribution, Significance, Originality, Overall Recommendation,
Confidence, and a constructive Comment as distinct values. A deterministic
validator checks required fields, integer ranges, non-empty prose, lease and
paper identity, evidence locations, and forbidden filler text. At most one
targeted repair may address a schema defect; it may not replace the frozen
paper or evidence.

\section{Reliability Model}

The coordinator uses a bounded queue with a default high-water mark of three,
one active lease per paper, append-only ledger writes with atomic snapshot
replacement, and idempotency keys for post intents. A timed-out worker may be
re-queued once under the bounded retry policy while Root retains the
deterministic lease. Malformed output records a validator-reject state and its
validation errors before it can enter the single targeted repair path.

Claim and post timeouts are ambiguous side effects. The coordinator therefore
queries status before issuing the same action again. A post attempt remains
\texttt{post\_unknown} until remote state resolves it to verified success or a
safe retry. The only live completion condition is
\texttt{posted\_verified == assigned\_count}. Mock completion is reported as
mock evidence and never as live platform performance.

\section{Evaluation Protocol}

Evaluation is split into quality and runtime contracts. Quality cases are
seeded with issue type and evidence location before any agent call. A naive
single-pass reviewer is frozen as the baseline. Issue matching reports true
positives, false positives, and false negatives; location matching follows the
predeclared page, section, table, and figure rules. The acceptance threshold is
stored with the fixture rather than chosen after observing results.

Runtime evaluation uses ten mock assignments and three independent repetitions.
It records assigned and verified-complete counts, schema validity, duplicate
post count, and latency from raw event records. Separate fault scenarios inject
malformed JSON, worker timeout, claim timeout, and post timeout, and reopen the
ledger once in-process. Each scenario must show either verified recovery or an
explicit unresolved state; silent loss is not a pass.

\section{Synthetic Mock Results}

Table~\ref{tab:runtime} reports only deterministic synthetic fixtures and the
mock adapter. Across three normal runs, all 30 assigned paper instances reached
mock \texttt{posted\_verified}; all 30 drafts passed the schema, and duplicate
post attempts and duplicate idempotency keys were both zero. A separate
10-paper run injected one malformed JSON output, worker timeout, claim timeout,
and post timeout, and reopened the ledger once in-process. Each injected
condition was recovered once, all ten drafts were valid and mock-verified, and
duplicate post attempts remained zero. A no-op rerun after full completion left
the ledger and outbox byte-identical.

Actual process recovery was tested separately. One process durably exited with
code 75 after four mock-verified papers. A fresh process opened the same output
directory and inputs, exited 0 at 10/10 mock verification with schema validity
10/10 and zero duplicate posts or keys, preserved the four-paper ledger prefix,
and preserved the completed manifest and outbox bytes.

\begin{table}[t]
  \caption{Deterministic synthetic-mock runtime evidence. No production claim
  or post occurred.}
  \label{tab:runtime}
  \centering
  \small
  \begin{tabular}{@{}lrrr@{}}
    \toprule
    Run & Verified & Schema & Dup. \\
    \midrule
    Normal, three repetitions & 30/30 & 30/30 & 0 \\
    Fault injection & 10/10 & 10/10 & 0 \\
    Two-process resume & 10/10 & 10/10 & 0 \\
    \bottomrule
  \end{tabular}
\end{table}

On the pre-frozen seeded synthetic issue set, the deterministic reviewer found
20 true positives, zero false positives, and zero false negatives, yielding F1
1.0. The naive single-pass baseline found 10 true positives, zero false
positives, and 10 false negatives, yielding F1 0.6667. Both had location
accuracy 1.0 under the frozen location rule. These results test deterministic
fixture recognition, not scientific judgment or generalization.

\begin{table}[t]
  \caption{Seeded synthetic-issue quality evaluation.}
  \label{tab:quality}
  \centering
  \small
  \begin{tabular}{@{}lrrrr@{}}
    \toprule
    Reviewer & TP & FP & FN & F1 \\
    \midrule
    Deterministic runtime & 20 & 0 & 0 & 1.0000 \\
    Naive single pass & 10 & 0 & 10 & 0.6667 \\
    \bottomrule
  \end{tabular}
\end{table}

Runtime and quality values come from the frozen
\nolinkurl{evidence/mock-validation-final-20260712T1335KST/aggregate.json}.
The actual process result comes from
\nolinkurl{evidence/process-restart-proof/aggregate.json}. Values above are
direct aggregates or four-decimal rounding of their stored values.

\section{Limitations}

Passing mock runs cannot establish live platform throughput, submission
success, review correctness on unseen research, or agreement with human
judgment. The runtime elapsed times reflect in-process synthetic fixtures and
are intentionally not used as a live speed claim. Worker-timeout evidence uses
a mock adapter \texttt{TimeoutError} and bounded recovery; it does not prove
that a running thread was forcibly killed.

Further limitations are structural. Evidence locations can show where a claim
was derived but cannot guarantee the claim is correct. Originality assessment
is limited to the submitted paper and its supplied references. A deterministic
schema validator can reject malformed reviews but cannot certify scientific
quality. The live adapter remains unspecified until read-only discovery reveals
the actual platform contract; no endpoint or selector is assumed here.

\section{Conclusion}

The design combines frozen evidence identity, a narrow side-effect owner,
bounded paper-local workers, deterministic validation, and status-first
reconciliation. Its evaluation contract separates review quality from runtime
reliability and separates mock evidence from live outcomes. Deterministic mock
runs satisfy the frozen runtime, seeded-fixture, and two-process resume checks,
while live platform behavior and unseen-paper review quality remain unmeasured.

\begin{center}
  \small\texttt{BODY-END-MARKER}
\end{center}

\begin{thebibliography}{2}

\bibitem[Ralphthon ICML organizers(2026a)]{officialskill}
Ralphthon ICML organizers.
\newblock Auto Research Skill, Track 2-only frozen-paper path.
\newblock Repository commit
\nolinkurl{a9f4f2583648ef4ca54f980f951ae393d153473f}, 2026.

\bibitem[Ralphthon ICML organizers(2026b)]{eventguide}
Ralphthon ICML organizers.
\newblock Track 2 rules, 2026.

\end{thebibliography}

\end{document}
